\answer{ex:21:1}{(a)~$T=\tau_{\text{ph}}\ln\frac{1}{1-x^*}$: $1.84$~s và $3.68$~s. (b)~Từ $3.36$ đến $4.24$~s: độ nhạy $\dd T/\dd x^*=\tau_{\text{ph}}/(1-x^*)=80$~s, khi $x^*\to1$ các đợt bùng phát không đều.}
\answer{ex:21:2}{(a)~$8.6$~W, như trong~\eqref{eq:energy}. (b)~$205$~Mbit/s, $0.2$~W, gấp $2\cdot10^3$ lần mẫu nuôi cấy ($10^{-4}$~W).}
\answer{ex:21:3}{$x_{\text{dừng}}=1/(1+Ur_b\tau_{\text{ph}})<x^*$ khi $r_b>(1/x^*-1)/(U\tau_{\text{ph}})=0.28$~Hz với $x^*=0.9$ và $0.025$~Hz với $x^*=0.99$; sức mạnh khớp giảm xuống $x_{\text{dừng}}$, tức giảm $10\%$ và $1\%$.}
\answer{ex:21:4}{Các $u(t-k)$ trực chuẩn, $m_k=\|P_Vu(t-k)\|^2$, trong đó $P_V$ là phép chiếu lên bao tuyến tính $N$ chiều của các đầu ra; theo bất đẳng thức Bessel, $\sum_k m_k\le N$.}
\answer{ex:21:5}{$\mathrm{MC}\approx9$--$11$ khi có $\tanh$ và $15$--$18$ với hồ chứa tuyến tính (ba hạt giống ngẫu nhiên), cả hai đều $\le30$: tính phi tuyến phải trả giá bằng trí nhớ.}
\answer{ex:21:6}{(a)~$n\le102$. (b)~$5.6$ sigma.}
\answer{ex:21:7}{Bài mở. Đối chứng then chốt: đóng băng lớp đọc ra và kiểm tra xem sự cải thiện có còn sau một khoảng nghỉ không kích thích hay không; nếu không, cái thay đổi là trạng thái, không phải trọng số.}
